% RESULTADO_RECUPERADO desde II REV09. Véase MAPA_PROCEDENCIA.json.
% Selección de líneas y cambios mecánicos explícitos; ninguna prueba nueva.
\section{Incidencia hexádico--octádica y funcional de Barbero--Immirzi}
\label{v:ant:sec:barbero}
Sea $H\hookrightarrow O$ una inclusión marcada, $|H|=6$, $|O|=8$.
Los objetos de incidencia
\[
\mathcal F_6=H\times\binom H2,\qquad
\mathcal F_8=O\times\binom H2
\]
tienen cardinales $90$ y $120$. Su diferencia normalizada da
\[
\frac{|\mathcal F_6|-|\mathcal F_8|}{24\binom62}=-\frac1{12}.
\]
\begin{figure}[htbp]
\centering
\begin{tikzpicture}[>=Latex,every node/.style={font=\small}]
\draw[rounded corners,draw=ink] (-1.8,-1.35) rectangle (1.8,1.35);
\foreach \a in {0,60,...,300}{
\fill[ink] (\a:0.83) circle(2pt);}
\node at (0,-1.7) {$H:\ 6$ puntos};
\draw[->,thick,ink] (2.1,0)--(3.6,0);
\node[above] at (2.85,1.5) {inclusión marcada};
\draw[rounded corners,draw=ink] (4,-1.35) rectangle (7.6,1.35);
\foreach \a in {0,60,...,300}{
\fill[ink] ($(5.8,0)+(\a:0.83)$) circle(2pt);}
\fill[accent] (4.5,0) circle(2pt);
\fill[accent] (7.1,0) circle(2pt);
\node at (5.8,-1.7) {$O:\ 8$ puntos};
\node at (0,-2.35) {$6\binom62=90$};
\node at (5.8,-2.35) {$8\binom62=120$};
\end{tikzpicture}
\caption{Incidencia marcada de seis a ocho puntos. El dibujo representa
los conjuntos y la inclusión; no pretende convertir las caras o vértices
de un cubo en bloques de Steiner sin su mapa de incidencia.}
\end{figure}

\subsection{Los pesos preceden al coeficiente de polo}
En el módulo libre de sucesos orientados, la fuente fija la cadena
fundamental de la vuelta
\[
c_{12}^{+}=
\sum_{j\in\Z/12\Z}[j,\mathcal F_6]
-[\partial_{\rm ret},\mathcal F_8].
\]
La multiplicidad doce corresponde a los sectores y la multiplicidad uno
a la costura orientada. La involución cambia la orientación de la cadena
completa. Sea
\[
S_m(q)=\frac{q^m}{1-q^{3m}},\qquad 0<q<1.
\]
El lector lineal asigna $S_{90}$ a cada sector y $S_{120}$ a la costura;
en consecuencia,
\[
\mathscr R(c_{12}^{+})=12S_{90}-S_{120}.
\]

\begin{theorem}[Coeficiente de polo y evaluación conjugada]
La imagen de la cadena fundamental tiene coeficiente $1/24$
respecto de $(1-q)^{-1}$ y residuo $-1/24$ respecto de $(q-1)^{-1}$.
Su evaluación en los canales previamente generados determina
\begin{equation}
\gamma_{\HMT}=
\frac{\pi AC^*}{180}
+12\bigl[S_{90}(q_-)-S_{90}(q_+)\bigr]
-\bigl[S_{120}(q_-)-S_{120}(q_+)\bigr].
\label{v:ant:eq:gamma}
\end{equation}
\end{theorem}
\begin{proof}
La factorización $1-q^{3m}=(1-q)\sum_{j=0}^{3m-1}q^j$ implica
$\lim_{q\uparrow1}(1-q)S_m(q)=1/(3m)$. Por linealidad,
\[
\frac{12}{270}-\frac1{360}=\frac1{24}.
\]
La coordenada $q-1$ invierte el signo del residuo.
La evaluación conjugada y la componente angular de la fuente
producen \eqref{v:ant:eq:gamma}.
\end{proof}

El certificado diofántico $4a-3b=45$ admite los pares
$(12+3t,1+4t)$, $t\ge0$; certifica la minimalidad de $(12,1)$,
pero no reemplaza el origen de las multiplicidades en la cadena.
La evaluación de \eqref{v:ant:eq:gamma} es
\[
\gamma_{\HMT}
=0.27400767269445927474725526077398041940\ldots.
\]
Esta cifra es una salida del funcional fuente; no es un valor externo
usado para seleccionar $A$, $C^*$ o sus coeficientes.

\subsection{Paridad y sensibilidad estructural}
Al invertir $C^*$ se intercambian $q_+$ y $q_-$ y cambia de signo
la componente angular. Por ello
\[
\gamma_{\HMT}(A,-C^*)=-\gamma_{\HMT}(A,C^*).
\]
El funcional conserva orientación, mientras que otros lectores,
como $q_+q_-$, son pares. Esta diferencia precede a su interpretación física.
